% Praca magisterska - Michał Cichoń, AGH 2014
% Title: Silnik do automatycznej kategoryzacji obrazów
\documentclass[a4paper,12pt]{book}
\usepackage[utf8]{inputenc}
\usepackage[T1]{polski}
\usepackage{polski}
\usepackage{color}
\usepackage{helvet}
\usepackage{graphicx}
\usepackage{geometry}
\usepackage{epstopdf}
\usepackage{titlesec}
\usepackage{indentfirst}
\usepackage{verbatim}
\usepackage{moreverb}
\usepackage{nameref}
\usepackage{listings}
\usepackage{url}
\usepackage[font=footnotesize,labelfont=bf]{caption}
\geometry{hmargin={2cm, 2cm}, height=10.0in}
\assignpagestyle{\chapter}{empty}
\linespread{1.3} %interlinia ustawiona na 1.5; 1.3 (LaTeX) to 1.5, 1.6 (LaTeX) to 2

\newenvironment{dedication}
{
   \cleardoublepage
   \thispagestyle{empty}
   \vspace*{\stretch{10}}
   \hfill\begin{minipage}[t]{0.66\textwidth}
   \raggedright
}%
{
   \end{minipage}
   \vspace*{\stretch{3}}
   \clearpage
}



\begin{document}
\nocite{*}

% =====  STRONA TYTULOWA PRACY MAGISTERSKIEJKIEJ ====
% ostatnia modyfikacja: 2009/07/01, K. Malarz

\thispagestyle{empty}
%% ------------------------ NAGLOWEK STRONY ---------------------------------
\includegraphics[height=37.5mm]{agh_nzw_a_pl_1w_wbr_cmyk.eps}\\
\rule{30mm}{0pt}
{\large \textsf{Wydział Fizyki i Informatyki Stosowanej}}\\
\rule{\textwidth}{3pt}\\
\rule[2ex]
{\textwidth}{1pt}\\
\vspace{7ex}
\begin{center}
{\LARGE \bf \textsf{Praca magisterska}}\\
\vspace{13ex}
% --------------------------- IMIE I NAZWISKO -------------------------------
{\bf \Large \textsf{Michał Cichoń}}\\
\vspace{3ex}
{\sf\small kierunek studiów:} {\bf\small \textsf{informatyka stosowana}}\\
\vspace{1.5ex}
{\sf\small specjalność:} {\bf\small \textsf{grafika komputerowa i przetwarzanie obrazów}}\\
\vspace{10ex}
%% ------------------------ TYTUL PRACY --------------------------------------
{\bf \huge \textsf{Silnik do automatycznej kategoryzacji obrazów}}\\
\vspace{14ex}
%% ------------------------ OPIEKUN PRACY ------------------------------------
{\Large Opiekun: \bf \textsf{dr inż.\ Maciej Śniechowski}}\\
\vspace{22ex}
{\large \bf \textsf{Kraków, wrzesień 2014}}
\end{center}
%% =====  STRONA TYTUŁOWA PRACY MAGISTERSKIEJKIEJ ====

\newpage

%% =====  TYŁ STRONY TYTUŁOWEJ PRACY MAGISTERSKIEJKIEJ ====
{\sf Oświadczam, świadomy odpowiedzialności karnej za poświadczenie nieprawdy, że niniejszą pracę dyplomową wykonałem osobiście i samodzielnie i  nie korzystałem ze źródeł innych niż wymienione w pracy.}

\vspace{14ex}

\begin{center}
\begin{tabular}{lr}
~~~~~~~~~~~~~~~~~~~~~~~~~~~~~~~~~~~~~~~~~~~~~~~~~~~~~~~~~~~~~~~~~ &
................................................................. \\
~ & {\sf (czytelny podpis)}\\
\end{tabular}
\end{center}

%% =====  TYL STRONY TYTULOWEJ PRACY MAGISTERSKIEJKIEJ ====

\newpage
\rightline{Kraków, ?? czerwca 2014}
\begin{center}
{\bf Tematyka pracy magisterskiej i praktyki dyplomowej
Michała Cichonia,
studenta V roku studiów kierunku informatyka stosowana, specjalności grafika komputerowa i przetwarzanie obrazów.}\\
\end{center}

Temat pracy magisterskiej:
{\bf Silnik do automatycznej kategoryzacji obrazów}\\

\begin{tabular}{rl}

Opiekun pracy:                  & dr inż.\ Maciej Śniechowski\\
Recenzenci pracy:               & ...\\
Miejsce praktyki dyplomowej:    & WFiIS AGH, Kraków\\
\end{tabular}

\begin{center}
{\bf Program pracy magisterskiej i praktyki dyplomowej}
\end{center}

\begin{enumerate}
\item Omówienie realizacji pracy magisterskiej z opiekunem.
\item Zebranie i opracowanie literatury dotyczącej tematu pracy.
\item Praktyka dyplomowa:
\begin{itemize}
\item zapoznanie się z ideą...,
\item uczestnictwo w eksperymentach/przygotowanie oprogramowania...,
\item dyskusja i analiza wyników...
\item sporządzenie sprawozdania z praktyki.
\end{itemize}
\item Kontynuacja obliczeń związanych z tematem pracy magisterskiej.
\item Zebranie i opracowanie wyników obliczeń.
\item Analiza wyników obliczeń numerycznych, ich omówienie i zatwierdzenie przez opiekuna.
\item Opracowanie redakcyjne pracy.
\end{enumerate}


\noindent
Termin oddania w dziekanacie: ?? września 2014\\[1cm]

\begin{center}
\begin{tabular}{lcr}
.............................................................. & ~~~ &
.............................................................. \\
(podpis kierownika katedry) & & (podpis opiekuna) \\
\end{tabular}
\end{center}

\newpage

\noindent
Na kolejnych dwóch stronach proszę dołączyć kolejno recenzje pracy popełnione przez Opiekuna oraz Recenzenta (wydrukowane z systemu MISIO i podpisane przez odpowiednio Opiekuna i Recenzenta pracy). Papierową wersję pracy (zawierającą podpisane recenzje) proszę złożyć w dziekanacie celem rejestracji co najmniej na tydzień przed planowaną obroną.

\newpage

\noindent
Na kolejnych dwóch stronach proszę dołączyć kolejno recenzje pracy popełnione przez Opiekuna oraz Recenzenta (wydrukowane z systemu MISIO i podpisane przez odpowiednio Opiekuna i Recenzenta pracy). Papierową wersję pracy (zawierającą podpisane recenzje) proszę złożyć w dziekanacie celem rejestracji co najmniej na tydzień przed planowaną obroną.

\begin{dedication}
Chciałbym złożyć serdeczne podziękowania Panu dr.~Maciejowi Śniechowskiemu za okazaną pomoc, przy realizacji tej pracy magisterskiej. A także moim rodzicom...
\end{dedication}

\tableofcontents

\chapter*{Wstęp}

%TODO TEN WSTEP WYMAGA POPRAWEK
\addcontentsline{toc}{chapter}{Wstęp}
Jednym z największych wyzwań, przed którymi stoi współczesna informatyka, jest przetwarzanie dużych ilości danych. Coraz częściej dane te nie reprezentują informacji liczbowych lub tekstowych, ale zawierają pewien rodzaj obrazu. 

Dobrą ilustracją tego zjawiska jest przykład przeglądarki internetowej Google, gdzie początkowo kładziony był nacisk na wyszukiwanie danych tekstowych, a obecnie obrazy są bardzo dynamicznie rozwijanym zasobem, którego indeksowanie w latach 2001-2010 zwiększyło się z 250 mln do ponad 10 mld.\cite{Google2010} 

Obrazy cyfrowe mają szereg zastosowań, od bardzo prozaicznych takich jak reprezentacja amatorskich zdjęć wykonanych cyfrowym aparatem fotograficznym, do bardzo zaawansowanych takich jak chociażby diagnostyka medyczna, gdzie informacja zawarta na zdjęciu jest kluczowa i niekiedy decyduje o ludzkim życiu.

Zarządzając dużym zbiorem obrazów często posługujemy się opisem słownym. Opis taki, w~formie nazwy pliku, metadanych EXIF\footnote{EXIF - Exchangeable Image File Format, najpopularniejszy standard metadanych dla plików ze zdjęciami cyfrowymi} lub rekordu w bazie danych, daje nam pewne wyobrażenie o tym co znajduje się na obrazie. Wyobrażenie to może być jednak mylne, lub niedostateczne, by pozwalało nam na poprawne jego zidentyfikowanie.

Pomoc w rozwiązaniu tego problemu daje nam teoria rozpoznawania obrazów, stosunkowo nowa, ale bardzo szybko rozwijająca się gałąź informatyki, która daje nam narzędzia umożliwiające takie przetwarzanie obrazów by możliwe stało się rozpoznanie przynależności obrazu do pewnej klasy.\cite{Tad91}

%TODO NAPISAĆ O CRIB (Wyszukiwanie informacji w dużych zbiorach obrazów)

\section*{Cel pracy}
\addcontentsline{toc}{section}{Cel pracy}

%TODO KONIECZNIE TRZEBA TO ZMIENIĆ - wstęp nieproporcjonalnie za długi w stosunku
%TODO do faktycznego celu pracy, domeny (2) też są nienajlepsze
Stworzenie systemu informatycznego, który pomógłby skategoryzować obrazy na podstawie ich zawartości mogłoby okazać się bardzo pomocne w wielu dziedzinach gospodarki, w których znalezienie odpowiedniej ilustracji do prezentowanych danych jest kluczowe takich jak: biznes, marketing, medycyna.

Celem pracy jest stworzenie silnika kategoryzacyjnego, który na podstawie danych wyekstratowanych z obrazu będzie w stanie przyporządkować obraz do najbardziej odpowiedniej kategorii.

\section*{Przykłady implementacji}
\addcontentsline{toc}{section}{Przykłady implementacji}
Napisać o:

- http://imagga.com/solutions/categorization-api.html

- http://googleblog.blogspot.com/2010/07/ooh-ahh-google-images-presents-nicer.html (problem z tulipanem)

\addcontentsline{toc}{chapter}{Bibliografia}
%\bibliography{praca_magisterska}

%%%%%%%%%%% Bibliografia %%%%%%%%%%%%%%%%%%%%%%%%%% 
\begin{thebibliography}{100} % 100 is a random guess of the total number of 
%references 
 
\bibitem{Boney96} Boney, L., Tewfik, A.H., and Hamdy, K.N., ``Digital 
Watermarks for Audio Signals," \emph{Proceedings of the Third IEEE 
International Conference on Multimedia}, pp. 473-480, June 1996. 
 
\bibitem{MG} Goossens, M., Mittelbach, F., Samarin, \emph{A LaTeX 
Companion}, Addison-Wesley, Reading, MA, 1994. 
 
\bibitem{HK} Kopka, H., Daly P.W., \emph{A Guide to LaTeX}, 
Addison-Wesley, Reading, MA, 1999. 
 
\bibitem{Pan} Pan, D., ``A Tutorial on MPEG/Audio Compression," \emph{IEEE 
Multimedia}, Vol.2, pp.60-74, Summer 1998. 
 
\end{thebibliography} 
%%%%%%%%%%%%% end %%%%%%%%%%%%%%%%%%%%%%%%%%%%%%% 


\end{document}

